\chapter{Background}
\label{ch:background}

This chapter introduces the background needed for the rest of the thesis: universal hashing, polynomial hash functions, prime-field arithmetic, limb representations, and the SIMD layout used by the generated code.

\section{Universal and Delta-Universal Hashing}
\label{sec:universal-hashing}

Let \(\mathcal{R}\) be a key space, \(\mathcal{M}\) a message space, and \(\mathcal{D}\) an output domain. A keyed hash family is a function
\[
    H : \mathcal{R} \times \mathcal{M} \rightarrow \mathcal{D}.
\]
For a key \(r \in \mathcal{R}\), we write \(H_r(M)\) for \(H(r,M)\). A family is \(\varepsilon\)-almost universal if, for all distinct messages \(M,M' \in \mathcal{M}\),
\[
    \Pr_{r \leftarrow \mathcal{R}}\left[H_r(M) = H_r(M')\right] \leq \varepsilon.
\]
In many cryptographic applications one needs the stronger delta-universal variant. If \(\mathcal{D}\) is an additive group, \(H\) is \(\varepsilon\)-almost delta-universal if, for all \(M \neq M'\) and all \(Z \in \mathcal{D}\),
\[
    \Pr_{r \leftarrow \mathcal{R}}\left[H_r(M) - H_r(M') = Z\right] \leq \varepsilon.
\]
These definitions, originating in universal hashing~\cite{carter1979universal}, combine with pseudorandom masks or block-cipher calls to build message authentication codes. The construction inherits computational security from the masking primitive, while the collision behavior of the hash family is bounded information-theoretically.

\section{Polynomial Hash Functions}
\label{sec:polynomial-hash-functions}

Polynomial hash functions instantiate \(H\) by evaluating a polynomial over a finite field~\cite{bernstein2007pema}. The classical univariate construction interprets message blocks \(M_1,\ldots,M_n\) as field elements and evaluates
\[
    H_{\mathrm{poly}}(r,M)
        = \sum_{i=1}^{n} M_i r^{n-i}
        = M_1 r^{n-1} + M_2 r^{n-2} + \cdots + M_n.
\]
This form is convenient because it can be evaluated online by Horner's rule:
\[
    T \leftarrow M_1,\qquad
    T \leftarrow rT + M_i \quad \text{for } i=2,\ldots,n.
\]
The degree controls the universality bound, which also depends on the message domain and byte-to-field encoding. The encoding must be injective on the intended message space.

\subsection{Poly1305}
\label{subsec:poly1305}

Poly1305 is the best-known prime-field polynomial hash function~\cite{bernstein2005poly1305aes}. It works over the field \(2^{130}-5\). Message bytes are split into blocks and interpreted as field elements; a polynomial is then evaluated at a secret key \(r\). In the complete authenticator, the result is combined with a second secret value \(s\).

Poly1305 is both a deployed baseline and an example of the tension between design and implementation. Later work observes that some original choices, such as key clamping, no longer provide the same benefit in modern integer and vectorized implementations~\cite{degabriele2024sok}.

\section{Prime Fields and Crandall Primes}
\label{sec:prime-fields}

This thesis focuses on prime fields of the form
\[
    \mathbb{F}_p,\qquad p = 2^\pi - \theta,
\]
where \(\theta\) is small. Such primes are often called Crandall primes and are widely used in efficient prime-field hashing implementations~\cite{degabriele2024sok,pegolotti2026sphgen}. They are attractive for implementation because powers of \(2^\pi\) can be reduced cheaply:
\[
    x \cdot 2^\pi \equiv x \cdot \theta \pmod{2^\pi - \theta}.
\]
Thus terms above bit position \(\pi\) can be folded into the low limbs by multiplying by \(\theta\).

\paragraph{Byte encoding.}
At the API boundary, messages and keys are byte strings. The compiler maps consecutive chunks of bytes to field elements. For field parameter \(\pi\), the chunk size used here is \(\lfloor \pi/8 \rfloor\) bytes. Thus \(2^{116}-3\) consumes 14 bytes per element, while \(2^{226}-5\) consumes 28. Fractional-byte capacity is ignored; the next byte belongs to the next field element.

\section{Limb Representations}
\label{sec:limb-representations}

Machine instructions operate on fixed-width words, while the field elements considered here are typically larger than one word. We write \(\omega\) for the machine-word size, for example \(\omega=64\). Field elements are represented as limbs stored in \(\omega\)-bit words. If the regular limb size is \(\lambda\), then the base is \(B=2^\lambda\), with \(\lambda \leq \omega\). An integer \(a\) represented by \(\ell\) limbs has the form
\[
    a = \sum_{i=0}^{\ell-1} a_i 2^{i\lambda},
\]
where each \(a_i\) is stored in one word. The most significant limb may use fewer bits. We denote its size by
\[
    \lambda' = \pi - (\ell-1)\lambda,
\]
so canonical limbs satisfy \(0 \leq a_i < 2^\lambda\) for \(i<\ell-1\) and \(0 \leq a_{\ell-1} < 2^{\lambda'}\).

\paragraph{Saturation.}
A representation is saturated if \(\lambda=\omega\) for the regular limbs, and unsaturated if \(\lambda<\omega\). In unsaturated representations, the spare \(\omega-\lambda\) high bits temporarily hold carries, allowing additions and partial products to accumulate before carry propagation.

\paragraph{Balance.}
Balanced representations use the same number of value bits in all limbs. In unbalanced representations, the most significant limb is smaller. Since \(\pi\) is rarely a multiple of \(\lambda\), unbalanced representations arise naturally and often avoid increasing the limb count.

\paragraph{Parameter choice.}
The field parameters \(\pi,\theta\), word size \(\omega\), and arithmetic bounds determine \(\lambda\), \(\lambda'\), and \(\ell\). SPHGen gives bounds for intermediate limbs after operations such as schoolbook multiplication~\cite{pegolotti2026sphgen}. We use the same criterion: choose \(\lambda\) as large as possible, while keeping the configured arithmetic and carry schedule within machine-word bounds.

\section{Prime-Field Arithmetic}
\label{sec:prime-field-arithmetic}

Field addition adds corresponding limbs independently. The result may not be canonical, but still represents the correct field element while limb bounds are respected.

\paragraph{Multiplication.}
For multiplication, let
\[
    a=\sum_{i=0}^{\ell-1} a_i 2^{i\lambda},
    \qquad
    b=\sum_{j=0}^{\ell-1} b_j 2^{j\lambda}.
\]
The ordinary schoolbook product is
\[
    ab = \sum_{k=0}^{2\ell-2}
             \left(\sum_{i+j=k} a_i b_j\right)2^{k\lambda}.
\]
A direct product therefore has up to \(2\ell-1\) limbs. For a Crandall prime \(p=2^\pi-\theta\), high terms can be folded back using \(2^\pi \equiv \theta \pmod p\). When the top limb is unbalanced, powers above \(\pi\) also carry a factor \(2^{\lambda-\lambda'}\), which we denote as \(\kappa\). The generated schoolbook multiplication follows the same principle as the usual folded product:
\[
    c_k =
      \sum_{i+j=k} a_i b_j
      +
      \theta \cdot 2^{\lambda-\lambda'}
      \sum_{i+j=k+\ell} a_i b_j,
      \qquad 0 \leq k < \ell,
\]
with shifts and masks specialized to \(\lambda\), \(\lambda'\), and \(\theta\).

\paragraph{Squaring.}
Squaring is a special case of multiplication. Since \(a_i a_j=a_j a_i\), cross terms can be computed once and doubled. This reduces the number of limb multiplications and is useful for SQH.

\paragraph{Karatsuba.}
Karatsuba multiplication is an alternative to schoolbook multiplication~\cite{KarOfm62}. The compiler uses one Karatsuba level. Let
\[
    m=\left\lceil \ell/2 \right\rceil,\qquad X=2^{m\lambda},
\]
and split the two operands as
\[
    a=a_0+Xa_1,\qquad b=b_0+Xb_1,
\]
where \(a_0,b_0\) contain the lower \(m\) limbs and \(a_1,b_1\) the upper limbs. If \(\ell\) is odd, the upper halves are zero-extended when forming sums. The compiler computes
\[
    z_0=a_0b_0,\qquad
    z_2=a_1b_1,\qquad
    z_1=(a_0+a_1)(b_0+b_1)-z_0-z_2.
\]
The unreduced product is reconstructed as
\[
    ab = z_0 + Xz_1 + X^2z_2.
\]
After this reconstruction, terms whose limb index is at least \(\ell\) are folded back into the lower \(\ell\) limbs using the same Crandall factor as in the schoolbook case. Equivalently, if \(d_k\) denotes the \(k\)-th limb coefficient of \(z_0+Xz_1+X^2z_2\), the generated partially reduced limbs have the shape
\[
    c_k = d_k + \theta \cdot 2^{\lambda-\lambda'} d_{k+\ell},
    \qquad 0 \leq k < \ell,
\]
with absent coefficients interpreted as zero. This saves limb multiplications, but introduces additions, subtractions, temporaries, and often more register pressure.

\section{Carry Propagation and Limb Realignment}
\label{sec:carry-realignment}

Carry propagation restores limb bounds. For limbs \(a_i\), a carry round extracts
\[
    t_i = \left\lfloor a_i / 2^\lambda \right\rfloor
\]
from regular limbs, adds \(t_i\) to the next limb, and masks \(a_i\) to \(\lambda\) bits. For the most significant limb the extracted value is
\[
    t_{\ell-1} = \left\lfloor a_{\ell-1} / 2^{\lambda'} \right\rfloor,
\]
and the Crandall identity sends it back to the least significant limb as \(t_{\ell-1}\theta\). A full reduction additionally handles the case where the represented integer is still at least \(p=2^\pi-\theta\).

\paragraph{Realignment.}
The term limb realignment is used for the analogous process after multiplication, where intermediate products occupy wider temporary words. In the IR, both operations appear as carry instructions: synchronization points that restore the representation expected by later arithmetic.

\paragraph{Delay strategies.}
Carry placement is a performance choice. Immediate realignment carries as soon as possible. Partially delayed realignment carries at selected loop boundaries. Fully delayed realignment omits intermediate carries and relies on the final reduction, provided the arithmetic bounds allow it.

\section{SIMD Model for Polynomial Hashing}
\label{sec:simd-model}

The generated vectorized code uses SIMD lanes to process several field elements in parallel. The layout is limb-major: a vector register holds the same limb of several independent field elements.

\paragraph{Lane layout.}
Along the limb dimension, several registers represent one or more big integers. Along the lane dimension, each register contains independent field elements: two on the NEON target and four on AVX2. This follows previous generated polynomial-hashing implementations~\cite{pegolotti2026sphgen} and suits sums and parallel Horner branches.
